% from ega4-2-fr.tex
\oldpage[IV-2]{225}
\phantomsection
\addcontentsline{toc}{section}{Index des notations}
\section*{INDEX DES NOTATIONS}

\begingroup
\small
\setlength{\parindent}{0pt}
\setlength{\parskip}{2pt}

$\mathfrak S(X)$ ($X$ préschéma) : 2.2.15 et 4.8.1.

$\operatorname{Ass}(\mathscr F)$ ($\mathscr F$ $\mathcal O_X$-Module) : 3.1.1.

$\operatorname{deg.tr.}_K L$ ($L$ extension du corps $K$) : 4.1.

$\lambda_x(\mathscr F)$ ($\mathscr F$ $\mathcal O_X$-Module, $x$ point maximal
de $\operatorname{Supp}(\mathscr F)$) : 4.7.5.

$\Phi(\mathscr F)$ ($\mathscr F$ $\mathcal O_X$-Module) : 4.8.1.

$\operatorname{codim}(Y,X)$ ($Y$ partie quelconque d'un préschéma $X$) : 5.1.3.

$\dim(\mathscr F)$ ($\mathscr F$ $\mathcal O_X$-Module) : 5.1.12.

$\operatorname{coprof}(\mathscr F)$, $\operatorname{coprof}(X)$ ($X$
préschéma, $\mathscr F$ $\mathcal O_X$-Module) : 5.7.1.

$\operatorname{coprof}_A(M)$ ($A$ anneau noethérien, $M$ $A$-module de type
fini) : 5.7.12.

$\mathcal H^0_{X/Z}(\mathscr F)$ ($X$ préschéma, $Z$ partie de $X$ stable par
spécialisation, $\mathscr F$ $\mathcal O_X$-Module) : 5.9.1.

$\rho_{X/Z}$ ($X$ préschéma, $Z$ partie de $X$ stable par
spécialisation) : 5.9.7.

$\operatorname{prof}_T(\mathscr F)$ ($T$ partie d'un préschéma $X$, $\mathscr F$
$\mathcal O_X$-Module) : 5.10.1.

$Z^{(n)}(X)$ ($X$ préschéma, $n$ entier $\geq0$) : 5.10.13.

$A^{(1)}$ ($A$ anneau intègre noethérien) : 5.10.17 et 7.2.1.

$\operatorname{gr}_{\mathscr J}^{\bullet}(\mathscr F)$ ($\mathscr F$
$\mathcal O_X$-Module, $\mathscr J$ Idéal de $\mathcal O_X$) : 6.10.1.

$U_{S_n}(\mathscr F)$, $U_{C_n}(\mathscr F)$ ($\mathscr F$
$\mathcal O_X$-Module) : 6.11.2.

$\operatorname{CM}(\mathscr F)$ ($\mathscr F$ $\mathcal O_X$-Module) : 6.11.3.

$U_{S_n}(X)$, $U_{C_n}(X)$, $\operatorname{CM}(X)$ ($X$ préschéma) : 6.11.4.

$\operatorname{Reg}(X)$, $\operatorname{Sing}(X)$ ($X$ préschéma) : 6.12.1.

$U_{R_n}(X)$ ($X$ préschéma) : 6.12.9.

$\operatorname{Nor}(X)$ ($X$ préschéma) : 6.13.1.

$k\{T_1,\ldots,T_n\}$ : 7.4.8.

\endgroup

\clearpage
